%================================================
\section{Results and Discussion}

%================================================
\subsection{Ensemble}
%================================================
\begin{frame}{Ensemble Trajectory Comparison}
	\vspace{-0.2cm}
	\begin{columns}[T]

		\column{0.48\textwidth}

		\centering
		\textbf{Experiment}

		\vspace{0.15cm}
		\includegraphics[height=0.62\textheight]{for_2.3_droplets_all_droplets_trajectory_comparison.png}
		\column{0.48\textwidth}

		\centering
		\textbf{Simulation}

		\vspace{0.15cm}

		\includegraphics[height=0.62\textheight]{output_sim_droplet_6_tracking_all_droplets_trajectory_comparison.png}

	\end{columns}

	\vspace{0.25cm}

	\centering
	\small

	Persistent exploration and self-avoiding behaviour are observed in both datasets.
	The trajectory start and end positions are indicated by
	$\circ$ and $\blacksquare$, respectively.

\end{frame}

%================================================
\begin{frame}{Ensemble Mean-Squared Displacement}

	\begin{columns}[T]

		%------------------------------------------------
		\column{0.48\textwidth}

		\centering
		\textbf{Experiment}

		% \vspace{-0.05cm}

		\includegraphics[height=0.52\textheight]{for_2.3_droplets_ensemble_for_2.3_droplets_loglog_msd.png}

		\vspace{-0.1cm}

		{\small
			$\alpha = 1.88 \pm 0.00$
		}

		%------------------------------------------------
		\column{0.48\textwidth}

		\centering
		\textbf{Simulation}

		% \vspace{0.15cm}
		% \vspace{-0.15cm}

		\includegraphics[height=0.52\textheight]{output_sim_ensemble_droplet_6_tracking_loglog_msd.png}

		\vspace{-0.1cm}

		{\small
			$\alpha = 1.53 \pm 0.01$
		}

	\end{columns}

	\vspace{-0.1cm}

	\begin{alertblock}{Key Observation}
		Both datasets exhibit superdiffusive transport ($\alpha > 1$), indicating persistent motion and departure from normal diffusion.
	\end{alertblock}

\end{frame}


%================================================
\begin{frame}{Ensemble Velocity Autocorrelation Function}
	\vspace{-0.2cm}

	\begin{columns}[T]

		%------------------------------------------------
		\column{0.48\textwidth}

		\centering
		\textbf{Experiment}

		\vspace{0.15cm}

		\includegraphics[width=\linewidth]{for_2.3_droplets_ensemble_for_2.3_droplets_vacf.png}

		\vspace{0.1cm}

		{\small
			$\tau_p = 125.3 \pm 11.1~\mathrm{s}$
		}

		%------------------------------------------------
		\column{0.48\textwidth}

		\centering
		\textbf{Simulation}

		\vspace{0.15cm}

		\includegraphics[width=\linewidth]{output_sim_ensemble_droplet_6_tracking_vacf.png}

		\vspace{0.1cm}

		{\small
			$\tau_p = 0.031 \pm 0.004$
		}

	\end{columns}

	\vspace{0.2cm}

	\begin{alertblock}{Key Observation}
		Both VACFs decay over a finite timescale, indicating persistent velocity correlations and finite translational memory.
	\end{alertblock}

\end{frame}
%================================================
\begin{frame}{Ensemble Orientation Correlation Function}

	\begin{columns}[T]

		%------------------------------------------------
		\column{0.48\textwidth}

		\centering
		\textbf{Experiment}

		\vspace{0.15cm}

		\includegraphics[width=\linewidth]{for_2.3_droplets_ensemble_for_2.3_droplets_orientation_corr.png}

		\vspace{0.1cm}

		{\small
			$\tau_r = 59.8 \pm 0.8~\mathrm{s}$
		}

		%------------------------------------------------
		\column{0.48\textwidth}

		\centering
		\textbf{Simulation}

		\vspace{0.15cm}

		\includegraphics[width=\linewidth]{output_sim_ensemble_droplet_6_tracking_orientation_corr.png}

		\vspace{0.1cm}

		{\small
			$\tau_r = 0.00305 \pm 0.00054$
		}

		% If using scaled values instead:
		% {\small
		% $\tau_r = 0.00305 \pm 0.00054$
		%
		% $\tau_r^{(\mathrm{scaled})}
		% =
		% 1.41 \pm 0.25~\mathrm{s}$
		% }

	\end{columns}

	\vspace{0.2cm}

	\begin{alertblock}{Key Observation}
		The orientation correlations decay gradually in both experiment and simulation, indicating finite directional persistence and rotational memory.
	\end{alertblock}

\end{frame}
%================================================

%================================================
\begin{frame}{Summary of Ensemble Transport Parameters}

	\centering
	\scriptsize

	\begin{table}
		\centering
		\begin{tabular}{lcccccc}
			\toprule
			Dataset      &
			$N$          &
			Frames /
			Iterations   &
			$\alpha$     &
			$\tau_p$ (s) &
			$\tau_r$ (s)                         \\
			\midrule

			2.1
			             & 4
			             & 3648
			             & $1.8345 \pm 0.0019$
			             & $61.83 \pm 0.36$
			             & $13.49 \pm 0.34$
			\\

			2.2
			             & 6
			             & 3363
			             & $1.6482 \pm 0.0053$
			             & $17.14 \pm 0.56$
			             & $3.96 \pm 0.26$
			\\

			2.3
			             & 6
			             & 3951
			             & $1.8776 \pm 0.0012$
			             & $125.25 \pm 1.09$
			             & $59.78 \pm 0.78$
			\\

			2.4
			             & 3
			             & 5267
			             & $1.7025 \pm 0.0049$
			             & $28.40 \pm 0.22$
			             & $32.18 \pm 0.35$
			\\

			2\&4 $\mu$L
			             & 4
			             & 5526
			             & $1.5592 \pm 0.0066$
			             & $14.32 \pm 0.16$
			             & $10.83 \pm 0.20$
			\\

			\midrule

			Simulation
			             & 6
			             & 1000
			             & $1.5272 \pm 0.0088$
			             & $0.0310 \pm 0.0036$
			             & $0.00305 \pm 0.00054$
			\\

			Simulation (Scaled)
			             & 6
			             & 1000
			             & $1.5272 \pm 0.0088$
			             & $14.32 \pm 1.66$
			             & $1.41 \pm 0.25$
			\\

			\bottomrule
		\end{tabular}
	\end{table}

	\vspace{0.2cm}

	{\tiny
		Simulation persistence times were converted using
		$t_{\mathrm{scale}}
			=
			\tau_p^{\mathrm{exp}}
			/
			\tau_p^{\mathrm{sim}}
			\approx 462$.
	}

	\begin{alertblock}{Key Observation}
		The model reproduces the superdiffusive transport exponent
		($\alpha>1$) and exhibits finite velocity and directional persistence,
		consistent with memory-driven dynamics.
	\end{alertblock}

\end{frame}

\begin{frame}{Experiment--Simulation Comparison}

	\centering

	\begin{tabular}{lcc}
		\toprule
		Feature                     & Experiment   & Simulation   \\
		\midrule
		Self-avoidance              & $\checkmark$ & $\checkmark$ \\
		Superdiffusion ($\alpha>1$) & $\checkmark$ & $\checkmark$ \\
		VACF decay                  & $\checkmark$ & Partial      \\
		OCF decay                   & $\checkmark$ & Partial      \\
		Translational memory        & $\checkmark$ & Partial      \\
		Rotational memory           & $\checkmark$ & Partial      \\
		Non-Markovian behaviour     & $\checkmark$ & Qualitative  \\
		\bottomrule
	\end{tabular}

	\vspace{0.2cm}
	\begin{minipage}{0.9\linewidth}
		\footnotesize
		\textbf{Limitation:}
		The simulation reproduces the qualitative decay of the VACF and OCF,
		but the persistence times ($\tau_p,\tau_r$) differ substantially from
		the experimental values.
	\end{minipage}
	% \textbf{Conclusion:}
	\begin{alertblock}{Conclusion}
		The PDE--SDE model captures the key transport and memory signatures
		(self-avoidance, superdiffusion, and finite correlation decay),
		but does not quantitatively reproduce the experimental persistence times.
		% 	The PDE--SDE model reproduces the main transport and memory signatures observed experimentally.
	\end{alertblock}

\end{frame}

%================================================
% \subsection{Summary of Results}

% \begin{frame}{Summary of Results}

% 	\begin{itemize}
% 		\item The experimental trajectories show self-avoiding motion.
% 		\item The MSD indicates superdiffusive transport.
% 		\item The local MSD exponent reveals a change in transport regime with lag time.
% 		\item The VACF and OCF show finite velocity and directional persistence.
% 		\item The coupled PDE--SDE model reproduces the main qualitative trends seen in experiment.
% 	\end{itemize}

% 	\vspace{0.3cm}

% 	\begin{alertblock}{Main Conclusion}
% 		The simulation captures the observed self-avoidance, transport scaling, and finite memory timescales.
% 	\end{alertblock}

% \end{frame}
